\section{Objetivos}
\subsection{ Objetivo General}
Analizar el comportamiento y las características operativas de un contactor comercial mediante el estudio de su circuito magnético variable y la medición experimental de sus parámetros eléctricos, fundamentado en las normativas técnicas nacionales e internacionales vigentes.

\subsection{ Objetivos Específicos}
\begin{itemize}
    \item Identificar mediante el desarme físico del contactor sus partes constitutivas esenciales, evaluando la función mecánica y magnética del núcleo, el martillo, las espiras de sombra y el entrehierro permanente.
    \item Determinar la variación de la reluctancia  e inductancia del circuito magnético del dispositivo en sus dos estados fundamentales de operación: reposo (desenergizado) y activado (energizado).
    \item Cuantificar experimentalmente la magnitud de la corriente de llamada en la bobina de control durante el transitorio de cierre, justificando el riesgo térmico de mantener este régimen de forma permanente.
    \item Medir las condiciones de régimen estable del dispositivo mediante la determinación de la corriente de mantenimiento y el establecimiento del umbral de la tensión mínima de caída o desenganche.
    \item Contrastar los valores eléctricos medidos en el banco de pruebas frente a las especificaciones técnicas suministradas por el fabricante en su catálogo.
\end{itemize}